% Propietario conservado: /Users/ruben/Documents/ChatGPT/jueces y controles/REVISION_ARGUMENTAL_APERTURAS_CIERRES_20260915/II/manuscript_es/sections/definicion_barbero_rev06.tex
\subsection{Definición estructural y función del parámetro de Barbero--Immirzi}
\label{ii:sec:ii-definicion-barbero}

\begin{definition}[Evaluación orientada de la incidencia hexádico--octádica]
El parámetro \(\gamma_{\HMT}\) es la evaluación conjugada de la
cadena fundamental \(c_{12}^{+}\), completada por su componente
angular. En el dominio \(0<q_\pm<1\), su definición
\eqref{ii:eq:gamma} se expresa como
\begin{equation}
 \gamma_{\HMT}
 =\frac{\pi AC^*}{180}
  +[\mathscr R(c_{12}^{+})](q_-)
  -[\mathscr R(c_{12}^{+})](q_+).
 \label{ii:eq:ii-definicion-barbero-evaluacion}
\end{equation}
La incidencia marcada proporciona los espacios de sucesos; la cadena
orientada fija sus multiplicidades; los canales de
\eqref{ii:eq:ii-canales} realizan la evaluación.
\end{definition}

La determinación utiliza conjuntamente los cardinales \(90,120\),
las doce contribuciones de sector y la costura de retorno. El
coeficiente de polo \(1/24\) y la minimalidad diofántica comprueban
la compatibilidad de esos pesos con la normalización. Su procedencia
se conserva en \(c_{12}^{+}\), antes de evaluar los canales. La
inversión angular intercambia \(q_+\) y \(q_-\) y produce
\(\gamma_{\HMT}(A,-C^*)=-\gamma_{\HMT}(A,C^*)\): el valor
adimensional transporta una orientación de la incidencia.

En la realización geométrica, ese mismo coeficiente interviene en
\eqref{ii:eq:holst} y en el incremento de área
\(\gamma_{\HMT}a_0\ell_P^2\). La acción de Holst utiliza
\(\gamma_{\HMT}\ne0\); la magnitud areal positiva emplea
\(|\gamma_{\HMT}|\). La proposición~\ref{ii:prop:ii-costura-barbero}
explicita el cambio coordinado de orientación. Así, la estructura
del parámetro explica qué pondera en la incidencia, qué conserva
al transportarse y cómo participa en la respuesta torsional y en
la resolución geométrica de la información.
